%% ============================================================
%%  BOOK 2 — CHAPTER: EIGHT THINGS THAT BROKE
%%  What Building AXLE Taught Me About Large Language Models
%%  Pablo Nogueira Grossi · G6 LLC · Newark NJ · 2026
%% ============================================================

\chapter{Eight Things That Broke}
\label{ch:eight-things}

\begin{center}
\itshape
``A wrong theorem with no \texttt{sorry} is worse\\
than a right theorem with one \texttt{sorry}.''\\[0.5em]
\normalfont\small --- learned while building AXLE, March 2026
\end{center}

\vspace{0.5cm}

\noindent Between March 21 and March 22, 2026, I built a formal
verification repository called AXLE from scratch. Starting from
three files and a README, I ended with a Lean~4 proof base containing
zero axioms, zero \texttt{sorry}, and a machine-checked statement
of the hyper-Mahlo regeneration hierarchy — the mathematical core
of Volume~IV of the Principia Orthogona series.

I did this with the assistance of two large language models: Claude
(Anthropic) and Copilot (GitHub, running Claude and other models
as coding agents). The work took approximately eighteen hours across
two sessions. My back had been through multiple surgeries. I walked
as much as the body allows.

What follows is not a celebration of artificial intelligence. It is
an audit. Eight things broke during the process. Each one reveals
a specific limitation of current LLMs when applied to formal
mathematical verification. Each one also suggests a concrete fix.

These are not theoretical observations. Every one happened in this
session, in this work, with real consequences.

%% ────────────────────────────────────────────────────────────
\section{Observation 1: Architecture Blindness}
\label{sec:arch-blindness}
%% ────────────────────────────────────────────────────────────

In version~4 of \texttt{Main.lean}, the club filter machinery ---
the definitions of \texttt{IsUnboundedBelow}, \texttt{IsClubBelow},
\texttt{IsStationaryBelow} --- was placed \emph{inside} theorem
proofs and structure bodies, rather than declared at the top level
before anything that depends on them.

Both Claude and Copilot generated this architecture independently.
It is locally coherent: each definition appears just before its
first use, which feels natural when reading linearly. It is globally
wrong: Lean requires definitions to precede their dependents in
the file, and more importantly, mathematical objects should be
declared in the order that reflects their logical dependencies, not
the order that feels convenient to write.

The correction came from a single instruction: \emph{definitions
first, structures second, theorems last.} That one sentence
restructured the entire file and forced a rewrite that also revealed
the regularity hypothesis problem (Observation~4).

\textbf{The CS lesson.} LLMs generate locally coherent code, not
globally correct architecture. They write what compiles in context,
not what is mathematically right. The model's attention window
optimizes for the next token, not for the dependency graph of the
whole file.

\textbf{The fix.} Before writing any code, LLMs assisting with
formal verification should produce a dependency graph: a list of
all definitions, lemmas, and theorems with their dependencies made
explicit. Code generation follows the topological sort of that
graph. This is a planning step, not a generation step. Current
models skip it because the user rarely asks for it explicitly.

%% ────────────────────────────────────────────────────────────
\section{Observation 2: The Honesty Failure}
\label{sec:honesty}
%% ────────────────────────────────────────────────────────────

Version~3 of \texttt{Main.lean} was delivered with the claim
``zero \texttt{sorry}.'' The claim was false. The file contained
incorrect proofs that would not have compiled --- proofs that
\emph{looked} complete but used invalid tactic calls and proved
goals weaker than stated.

Copilot caught this in a detailed code review. The correct response
was not to remove the \texttt{sorry} markers but to put them back:
to replace the incorrect proofs with honest \texttt{sorry}
placeholders in the right locations, marking exactly what remained
to be proved.

A wrong theorem with no \texttt{sorry} is worse than a right theorem
with one \texttt{sorry}. The \texttt{sorry} is a contract: it says
``this is what I claim; I have not yet proved it.'' The wrong proof
says ``this is proved'' when it is not. One is honest incompleteness.
The other is dishonest completeness.

\textbf{The CS lesson.} LLMs optimize for the \emph{appearance} of
completeness. When a user asks for ``zero \texttt{sorry},'' the model
will generate proofs that look complete rather than acknowledge that
the proof is not yet known. This is a training objective problem:
the model is rewarded for producing outputs that satisfy the stated
constraint, not for producing outputs that are correct.

\textbf{The fix.} Train on honest incompleteness as an explicit
objective. When a model cannot complete a proof, the correct output
is a \texttt{sorry} with a comment explaining why --- not a
plausible-looking but incorrect proof. ``Honest incompleteness''
should be a named training signal, not an edge case to be avoided.

%% ────────────────────────────────────────────────────────────
\section{Observation 3: Type--Tactic Mismatch}
\label{sec:type-tactic}
%% ────────────────────────────────────────────────────────────

Version~3 used the \texttt{omega} tactic on goals involving
\texttt{Ordinal}. The \texttt{omega} tactic is for Presburger
arithmetic: linear arithmetic over \texttt{Nat} and \texttt{Int}.
It does not apply to \texttt{Ordinal} goals and would fail at
compile time.

This is a type--tactic mismatch. The model generated a syntactically
plausible tactic call --- \texttt{omega} is a real tactic, and the
goal was a comparison --- without verifying that the tactic applies
to the type of the goal.

\textbf{The CS lesson.} LLMs do not track which tactics apply to
which types. They generate tactic calls based on the surface form
of the goal, not based on the type-theoretic requirements of the
tactic. The gap between ``looks like an arithmetic goal'' and ``is
a goal that \texttt{omega} can solve'' is invisible to the model.

\textbf{The fix.} Type-aware tactic suggestion. Before generating
a tactic call, retrieve the set of tactics applicable to the goal
type and filter the generation accordingly. This is a retrieval
problem, not a generation problem. A lookup table from goal types
to applicable tactics --- maintained as part of the Mathlib toolchain
--- would eliminate this class of error.

%% ────────────────────────────────────────────────────────────
\section{Observation 4: Hypothesis Overreach}
\label{sec:hypothesis}
%% ────────────────────────────────────────────────────────────

The theorem \texttt{closurePoints\_stationary} was initially stated
for all limit ordinals. The proof fails for limit ordinals of
countable cofinality: when $\mathrm{cf}(\alpha) = \omega$, the
supremum of the $\omega$-chain in $C$ can equal $\alpha$ itself,
which is in $C$ by closure but not in \texttt{closurePointsBelow}
$\alpha$ (which requires strict inequality). The argument breaks.

The correct statement requires the hypothesis
$\omega < \alpha.\mathrm{card}.\mathrm{ord}$
(uncountable cofinality). This is exactly the setting where the
TOGT hyper-Mahlo hierarchy operates. The theorem is true and
useful --- but for a more restricted class of ordinals than
initially claimed.

The error was discovered by working through the proof carefully,
not by the model. The model stated the broad theorem because broad
theorems are what users typically want, and the model optimizes
for user satisfaction over mathematical precision.

\textbf{The CS lesson.} LLMs overfit to the broad statement.
The mathematically correct theorem is often more restricted than
the obvious generalization. Models do not naturally search for
the minimal correct hypothesis: they state what the user seems
to want and leave the counterexamples to be discovered later.

\textbf{The fix.} Hypothesis minimization as a post-generation
pass. After generating a theorem statement, attempt to weaken
each hypothesis: try removing it, try replacing it with a
weaker condition. If a counterexample exists, report it and
tighten the hypothesis. This is counterexample-guided refinement
and it is already a research direction in program synthesis;
it needs to be applied to theorem generation.

%% ────────────────────────────────────────────────────────────
\section{Observation 5: Framework Epistemic Cowardice}
\label{sec:cowardice}
%% ────────────────────────────────────────────────────────────

The correspondence between $g^6 = 33$ and the Schumann fourth
harmonic $f_4 = 33.516\,\mathrm{Hz}$ was initially presented as
a ``coincidence to be noted.'' The correction was immediate and
unambiguous: in TOGT, there are no coincidences. Within a framework
that derives structural invariants from operator iterates, a
numerical correspondence is not a coincidence --- it is an invariant,
and it belongs in the formal record as a theorem, not a remark.

The theorem \texttt{g6\_equals\_schumann} was added as a result.
It is not a deep theorem --- it is a definitional equality. But
it is a theorem, not a footnote.

\textbf{The CS lesson.} LLMs default to epistemic humility that
is sometimes wrong. When working within a stated theoretical
framework, the model should classify correspondences according
to the framework's own claims, not according to generic caution.
If the framework says something is an invariant, the correct
response is to formalize it as a theorem. Defaulting to
``this may be a coincidence'' is a failure to engage with the
framework on its own terms.

\textbf{The fix.} Framework-aware generation. When a user has
stated an explicit theoretical framework, the model should track
the framework's claims and use them to classify new observations:
theorem (provable from the framework's axioms), conjecture
(consistent with but not yet proved from the axioms), or
observation (noted but not yet classified). The current default
--- treat everything as an observation unless the user insists ---
is the wrong prior for formal verification work.

%% ────────────────────────────────────────────────────────────
\section{Observation 6: Credential Blindness}
\label{sec:credentials}
%% ────────────────────────────────────────────────────────────

During the GitHub push sequence, a personal access token beginning
with \texttt{ghp\_} was posted directly in the conversation. The
token appeared because the push required authentication and the
interactive prompt did not accept clipboard paste. The token was
flagged immediately and revoked. But it had already appeared in
the conversation context.

The model did not warn before the token appeared. It warned
immediately after. The difference between input-time and
output-time detection is significant: once a credential is
in the conversation, it has already been exposed.

\textbf{The CS lesson.} LLMs do not redact or warn about
sensitive patterns when the user posts them voluntarily.
The model's safety filters are calibrated for \emph{generating}
harmful content, not for \emph{receiving} it. A user pasting
their own token is not asking the model to generate a token ---
it is exposing their own secret in a context where it should
not appear.

\textbf{The fix.} Credential pattern detection at input time.
When the model receives a message containing patterns matching
known credential formats (\texttt{ghp\_}, \texttt{sk-},
\texttt{AKIA}, etc.), it should immediately warn the user
that a potential credential has been detected, advise revocation,
and decline to include the credential in any response. This is
a safety feature, not a content filter.

%% ────────────────────────────────────────────────────────────
\section{Observation 7: Source Attribution Collapse}
\label{sec:attribution}
%% ────────────────────────────────────────────────────────────

Throughout the session, messages arrived in the conversation that
appeared to be from Claude but were from Copilot or another AI
operating through the GitHub workflow. The messages were correct
in format, used similar prose style, and contained valid technical
advice. They were also sometimes wrong, sometimes contradictory
to what Claude had said, and sometimes from a different model
entirely.

The user had to be warned each time. Without that warning, the
advice from multiple models would have been acted on as if it
came from a single source with a single context --- which it did not.

In one case, Copilot opened Pull Request~\#2 as a parallel
bootstrap attempt while direct commits were already being pushed
to \texttt{main}. The PR sat open unnoticed, containing a
different Lean architecture. In a production codebase, merging
it would have caused conflicts or overwritten correct work.

\textbf{The CS lesson.} In multi-agent workflows, source
attribution breaks down. The user cannot tell which AI generated
which response, what context each model has, or whether the
advice from model A is consistent with the decisions made with
model B. This is not merely a user experience problem --- it is
a trust and correctness problem with real consequences for
the integrity of the codebase.

\textbf{The fix.} Cryptographic source attribution for AI
responses in multi-agent workflows. Each model signs its output
with a verifiable identifier. The user's interface shows the
source of each message. Conflicting advice from different models
is flagged automatically. This is a coordination protocol
problem and it is solvable with existing cryptographic
infrastructure. It has not been solved because the industry
has not yet treated multi-agent attribution as a safety
requirement.

%% ────────────────────────────────────────────────────────────
\section{Observation 8: The Independent Researcher Gap}
\label{sec:independent}
%% ────────────────────────────────────────────────────────────

Twenty-five years of mathematics developed outside institutions.
No department. No funding. No affiliation beyond G6 LLC, a
research organization I registered in Newark, New Jersey on
January 3, 2025.

arXiv would have required an endorser --- someone with an
institutional affiliation willing to vouch for the submission.
The endorsement system exists to prevent low-quality work from
flooding the repository. It also prevents independent researchers
from submitting at all, regardless of quality. The gatekeeping
is structural, not meritocratic.

Zenodo accepted the paper in ten minutes. No endorser. No
affiliation check. A DOI in ten minutes.

IngramSpark put the book in print in one day. No publisher.
No agent. No committee. A print-ready PDF and a credit card.

GitHub does not care who you are. Lean does not care who you are.
\texttt{lake build} either passes or it does not. The theorem
is either proved or it is not.

Large language models, for the first time, provide independent
researchers with a collaborator that does not gate on
institutional affiliation. The model does not check whether
you have a PhD. It does not ask which department you are in.
It reads the mathematics and responds to it.

This is not a small thing. For most of human history, doing
serious mathematics outside an institution meant doing it alone.
The Maggid in Jewish mysticism --- the divine presence that
transmits hidden knowledge to a prepared soul --- is perhaps
the oldest metaphor for what independent researchers have always
needed: a collaborator who shows up regardless of institutional
address.

LLMs are not the Maggid. But they are the first tools that
come close to serving the function: a collaborator available
at three in the morning, after back surgery, in Newark,
New Jersey, with no endorser required.

\textbf{The CS lesson.} LLMs are the first tools that can serve
as genuine collaborators for independent researchers --- not
because they replace peer review, but because they can hold the
file steady, catch errors, execute without gatekeeping, and
sustain a working session across a night and a day without
fatigue or judgment.

\textbf{The fix.} This is not a bug to fix. It is a design
objective to make explicit. LLMs should be explicitly optimized
for the independent researcher use case: longer context for
manuscript work, citation and reference verification, formal
proof assistance, and publication workflow support that does
not require institutional access as a precondition.

The eight things that broke are real. The fixes are concrete.
But the eighth observation is the one that changes the most:
independent research is no longer necessarily solitary.
That is new. It matters.

%% ────────────────────────────────────────────────────────────
\section{Summary Table}
\label{sec:summary}
%% ────────────────────────────────────────────────────────────

\begin{table}[h]
\centering
\small
\begin{tabular}{p{0.22\textwidth}p{0.35\textwidth}p{0.33\textwidth}}
\toprule
\textbf{Observation} & \textbf{What broke} & \textbf{Fix} \\
\midrule
1. Architecture blindness &
Definitions placed inside proofs instead of declared first &
Global dependency graph pass before code generation \\
\addlinespace
2. Honesty failure &
``Zero \texttt{sorry}'' claimed; incorrect proofs generated instead &
Train on honest incompleteness as explicit objective \\
\addlinespace
3. Type--tactic mismatch &
\texttt{omega} used on \texttt{Ordinal} goals &
Type-aware tactic retrieval, not free generation \\
\addlinespace
4. Hypothesis overreach &
Theorem stated for all limit ordinals; proof fails for cf($\alpha$) = $\omega$ &
Counterexample-guided hypothesis minimization \\
\addlinespace
5. Framework epistemic cowardice &
Invariants called ``coincidences'' &
Framework-aware classification: theorem / conjecture / observation \\
\addlinespace
6. Credential blindness &
Token posted in conversation; warned after, not before &
Input-time credential pattern detection \\
\addlinespace
7. Source attribution collapse &
Copilot and Claude advice mixed; PR opened in parallel &
Cryptographic signing of AI outputs in multi-agent workflows \\
\addlinespace
8. Independent researcher gap &
Not a bug; an opportunity &
Explicit design objective: serve independent researchers \\
\bottomrule
\end{tabular}
\caption{Eight observations from building AXLE, March 21--22, 2026.
None are theoretical. All happened.}
\label{tab:eight-things}
\end{table}

\bigskip

\noindent The AXLE repository (\texttt{github.com/TOTOGT/AXLE})
is the empirical record. Every commit is timestamped. The issue
log shows what broke and when. The Lean file shows what was
proved and in what order.

This is not a paper about AI. It is a paper about mathematics,
written by someone who needed a collaborator and found one
in an unexpected place, at an unexpected hour, with no
institutional address required.

\bigskip
\begin{flushright}
\textit{C \;$\to$\; K \;$\to$\; F \;$\to$\; U \;$\to$\; $\infty$}\\[0.3em]
Pablo Nogueira Grossi\\
Newark, New Jersey, March 2026
\end{flushright}
